\section{Gaussian height windows and an explicit marked projection}\label{sec:gaussian-window}

We now change the physical window, while retaining the observation
Gaussian and the uniform center law. This change makes center integration
explicit and is useful for finite acquisition. Its tails must be paid;
the compact-window time constants of Section~\ref{sec:newman-time}
do not automatically apply to it. The analytical construction is from
\cite{OCGaussianWindow,OCProjection}.

\subsection{An exact center integral with controlled tails}

For $h>0$ use
\begin{equation}\label{eq:gaussian-height-window}
 W_X^{\rm G}(x+\ii y)=\ee^{-(x-X)^2/(2h^2)}\kappa_0(y),
 \qquad X\sim\operatorname{Unif}[T,2T].
\end{equation}
Here $h$ is a standard deviation, not the support radius of the earlier
compact window. For $m=(x+x')/2$ and $d=x-x'$,
\begin{equation}\label{eq:height-center-integral}
 \begin{gathered}
 \mathbb E_X\ee^{-(x-X)^2/(2h^2)}\ee^{-(x'-X)^2/(2h^2)}
   =I_{T,h}(m)\ee^{-d^2/(4h^2)},\\
 I_{T,h}(m)=\frac{h\sqrt\pi}{2T}
 \left[\operatorname{erf}\frac{2T-m}{h}
                    -\operatorname{erf}\frac{T-m}{h}\right].
 \end{gathered}
\end{equation}
The identity follows from
$(x-X)^2+(x'-X)^2=2(m-X)^2+d^2/2$ and a one-dimensional Gaussian integral.
Both finite center edges remain in $I_{T,h}$.

The full pair density has a similar exact center integration before
any projection. Put
$\sigma^2=h^2\epsilon^2/(h^2+\epsilon^2)$,
$\alpha=h^2/(h^2+\epsilon^2)$ and
$a_u(z)=x-\alpha\ell\Re u$. Completing the square gives
\[
 JW_X^{\rm G}(z)G_\epsilon(X+\ell u-z)
  =p_u(z)\ee^{-(X-a_u(z))^2/(2\sigma^2)},
\]
where
\[
 p_u(z)=\frac{J\kappa_0(y)}{2\pi\epsilon^2}
  \exp\left[-\frac{(y-\ell\Im u)^2}{2\epsilon^2}
             -\frac{\ell^2(\Re u)^2}{2(h^2+\epsilon^2)}\right].
\]
For each ordered distinct-label pair, multiply $p_u(z)p_v(w)$ by
the right side of \eqref{eq:height-center-integral} with $h=\sigma$,
$m=(a_u(z)+a_v(w))/2$, and $d=a_u(z)-a_v(w)$.
This simplifies integration over a supplied divisor; it does not
estimate the unknown divisor or its joint source field.

From \eqref{eq:time-input}, the total height weight is bounded by
$B_0(X)$, and the weight at $|x_j-X|\ge R$ by $B_R(X)$, where
\begin{equation}\label{eq:gaussian-weight-tail}
 B_R(X)=2A\sum_{n=R}^\infty
       [\log(2+X)+\log(n+2)]\ee^{-n^2/(2h^2)},\qquad R\in\mathbb Z_{\ge0}.
\end{equation}
Indeed each shell is covered by two unit intervals and their logarithms
are at most $\log(2+X)+\log(n+2)$. These bounds are uniform in
$\lambda\in[0,1/2]$. With $a_0=J/(2\pi\epsilon^2)$, keeping only
the fixed slab $[T-R,2T+R]$ loses at most
\begin{equation}\label{eq:gaussian-pair-tail}
 2a_0^2B_0(2T)B_R(2T)\quad\hbox{in the full pair},\qquad
 4a_0^2B_0(2T)B_R(2T)\quad\hbox{in the connected pair}.
\end{equation}
Every omitted pair has at least one omitted weight; connected subtraction
pays the product of means a second time.

\subsection{What linear Gaussian smoothing does to the source}

For the normalized real Gaussian $\gamma_h$ of variance $h^2$,
absolute integration of the complete cosine representation gives
\begin{equation}\label{eq:source-heat-shift}
 \int_\R\gamma_h(X-x)H_\lambda(x+\ii y)\dd x
       =H_{\lambda-h^2/2}(X+\ii y).
\end{equation}
Each frequency is multiplied by $\ee^{-h^2u^2/2}$.
This is a linear identity. It does not pass through the logarithm.
For $p_\lambda(z)=z^2-2\lambda$ at $\lambda=h^2/2$, the convolved
polynomial is $z^2$, whereas the original zeros are $\pm h$.
At the origin the Gaussian mean of $\log|x^2-h^2|$ is finite and
the log modulus of the convolved polynomial is $-\infty$.

The log-source formulas themselves remain valid for the noncompact
height window. To justify this, let
\[
 M_0=\max\left\{1,\int_0^\infty\ee^{u^2/2+11u}\Phi(u)\dd u\right\}.
\]
Then $|H_\lambda|\le M_0$ on every disk $D(q+4\ii,7)$, and
integration of \eqref{eq:time-input} along the safe line gives
$\log|H_\lambda(q+4\ii)|\ge\log H_0(4\ii)-A|q|\log(2+|q|)$.
Subharmonic circular means, integrated in radius, therefore imply
\begin{equation}\label{eq:noncompact-log-tail}
 \int_{[q,q+1]\times[-2,2]}|\log|H_\lambda||\dd^2z
 \le49\pi[2\log M_0-\log H_0(4\ii)+A|q|\log(2+|q|)].
\end{equation}
The rectangle fits inside the disk. The pointwise bound
$|\log|H||\le2\log M_0-\log|H|$ and the mean lower bound prove
the inequality, including the negative log near zeros.

The Laplacians of $f_{X,u}=JW_X^{\rm G}G_\epsilon(X+\ell u-\cdot)$
and $f_{X,u}f_{X,v}$ have Gaussian horizontal decay times a polynomial
of degree at most two; their transverse support stays in $[-2,2]$.
Multiply them first by a horizontal cutoff with uniformly bounded first
two derivatives. Formula \eqref{eq:noncompact-log-tail} makes the
cutoff error summable and vanishing, while
\eqref{eq:gaussian-weight-tail} handles the zero side. Passing to the
limit proves the absolutely convergent formulas
\begin{equation}\label{eq:gaussian-log-contraction}
 r_X(u)=\frac1{2\pi}\int\log|H_\lambda|\Delta f_{X,u},\qquad
 d_X(u,v)=\frac1{2\pi}\int\log|H_\lambda|\Delta(f_{X,u}f_{X,v}).
\end{equation}
Uniform absolute bounds justify products and the center average.
Their joint input is still $\mathbb E_X[L_X(\xi)L_X(\omega)]$,
with $L_X(\xi)=\log|H_\lambda(X+\xi)|$, or its covariance.
A whole-line comparison also needs the approximating field's log tail.
Neither \eqref{eq:source-heat-shift} nor the tail bounds acquire a
sharp estimate of this joint field.

\subsection{Integrate one marked observation test exactly}

Return to $\lambda=0$, keep $\epsilon/\ell=1/4$, and use the Gaussian
height window. The following bounded test preserves two separate
imaginary marks:
\[
 \Psi(u,v)=\exp\left[-\frac{(\Re u-\Re v)^2}2
                     -\frac{(\Im u-1/2)^2}2
                     -\frac{(\Im v+1/2)^2}2\right].
\]
Define the projection over the \emph{entire} observation plane by
\begin{equation}\label{eq:marked-projection}
 P_T=\mathbb E_X\int_{\C^2}\Psi(u,v)R_{2,X}(u,v)\dd^2u\dd^2v.
\end{equation}
The Jacobian $J=\ell^2$ makes each unweighted observation Gaussian
have mass one. A real difference has variance $2(1/4)^2=1/8$ and
each transverse observation has variance $1/16$. Using
\[
 \mathbb E\ee^{-(a+Z)^2/2}
   =(1+v)^{-1/2}\ee^{-a^2/[2(1+v)]},\qquad Z\sim N(0,v),
\]
which follows by completing the square, the three convolution factors
multiply to $c_G=32\sqrt2/51$. Thus
\begin{equation}\label{eq:projection-kernel-general}
 P_T=\sum_{j\ne k}K_{T,h,\ell}(z_j,z_k),\qquad
 K_{T,h,\ell}(z,w)=c_G I_{T,h}(m)
   \ee^{-[1/(4h^2)+4/(9\ell^2)]d^2}b_+(y)b_-(y'),
\end{equation}
where $b_\pm(y)=\exp[-\frac8{17}(y/\ell\mp1/2)^2]$.
Here $\kappa_0=1$ on the canonical strip. Tonelli's theorem applies
to the nonnegative sums and integrals, with finiteness supplied by
\eqref{eq:gaussian-weight-tail}.

For the finite acquisition below set $h=\ell=1$, $\epsilon=1/4$,
and write $I_T=I_{T,1}$, $K=K_{T,1,1}$. Then
\begin{equation}\label{eq:projection-kernel}
 K(z,w)=c_G I_T(m)\ee^{-25d^2/36}b_+(y)b_-(y'),\qquad
 D(z):=K(z,z)=c_G I_T(x)\ee^{-4/17}\ee^{-16y^2/17}.
\end{equation}
The ordering matters:
$K(z,w)/K(w,z)=\exp[\frac{16}{17}(y-y')]$.
The midpoint dependence and the two marks therefore cannot be replaced
by a function of the complex difference alone. This is a specified
scalar consumer of the full hierarchy, not a reconstruction of it.
